\phaseheader{8}{Add differentiable recovery-value guidance to the flow policy}{Attempt an algorithmic efficiency extension only after discrete critic reranking proves prospective actionability. Test whether a noise-step-conditioned critic can guide one action-generation trajectory as effectively as sampling and ranking four complete candidates.}

\subsection{Scientific hypothesis and relation to prior work}

Failure-informed or critic-guided diffusion/flow improvement is already an active area~\cite{afil,rover,qvgm,checkvla}. Therefore, ``using a critic to guide flow'' is not by itself a sufficient novelty claim. Within this project, guidance is valuable if it uses the causal recovery-value target in Equation~\eqref{eq:advantage}, operates only when the utility gate identifies a recoverable state, and achieves the same or better embodied recovery with lower matched cost.

Let $z_{\tau}$ be the noisy action latent and $v_{\theta}(z_{\tau},h_t,\tau)$ the frozen policy velocity. A generic guided form is
\begin{equation}
  v_{\mathrm{guided}}
  = v_{\theta}
  +s(\tau)\,\Pi\!(
    \nabla_{z_{\tau}}Q_{\phi}(h_t,z_{\tau},\tau)
  ),
  \label{eq:flow-guidance}
\end{equation}
where $s(\tau)$ is a frozen scale schedule and $\Pi$ clips or projects the gradient to preserve action bounds. The exact sign and parameterization must match the Xiaomi flow implementation; Equation~\eqref{eq:flow-guidance} is a design template, not evidence that the current detached client score is differentiable.

\subsection{What must be implemented}

\begin{enumerate}
  \item Expose the noisy action/token latent and flow integration step inside the Xiaomi model-side generation loop.
  \item Train or distill a critic conditioned on $z_{\tau}$ and $\tau$, using the same counterfactual recovery outcomes as Phase 5.
  \item Execute the critic inside the model server and preserve the original unguided action-only and feature-return behavior when guidance is disabled.
  \item Freeze guidance sign, scale schedule, integration-step window, suffix scope, gradient norm clipping, action deviation limit, and uncertainty fallback.
  \item Log intermediate critic values, gradient norms, clipping frequency, action deviation, NFE, backward passes, GPU time, and executed outcome.
\end{enumerate}

\implementationnote{Required architectural change.}{The current SAFE/best-of-$K$ scorer receives detached final candidate features on the client. True guidance requires a model-side gradient path through intermediate action latents. Implement this as an opt-in server protocol with focused tests and an unguided-equivalence smoke.}

\subsection{Development and locked test}

Develop on \code{ArrangeTea} and \code{CuttingToolSelection}, for example with 20 fresh identities per task. Tune one small guidance-scale grid and guided integration window on development identities only. Freeze the selected setting and test on the remaining current tasks plus held-out composite tasks from Phase 7.

Compare:
\begin{enumerate}
  \item unguided $K=1$;
  \item random $K=4$;
  \item critic-reranked $K=4$;
  \item critic-guided $K=1$;
  \item shuffled-critic or random-gradient negative control.
\end{enumerate}
Provide both matched-NFE and matched-wall-clock comparisons. An apparent gain obtained by spending more backward passes or latency than $K=4$ is not an efficiency result.

\subsection{Evaluation}

Primary outcomes are active-stage recovery and full-task success. Secondary outcomes are predicted versus realized $\Delta$, prior-stage regression, collision, saturation, oscillation, action-distribution shift, NFE, model calls, memory, GPU time, and end-to-end latency. Ablate guidance scale, guided step window, suffix length, terminal SAFE risk target versus recovery-value target, and the negative gradient control.

\positiveresult{Guided $K=1$ matches or exceeds critic-reranked $K=4$ within a frozen statistical margin, uses meaningfully less inference cost or latency, and does not increase collision, saturation, oscillation, or completed-stage regression. The frozen guidance setting transfers from the two development tasks to held-out tasks.}

\negativeresult{Retain critic reranking as the recovery executor. Do not present guidance as successful if it raises the critic score but not executed task outcomes, if a shuffled gradient performs similarly, or if matched-compute comparisons remove the benefit. Phase 8 is optional; the causal recoverability paper does not require it to pass.}

\subsection{Required outputs}

Release the opt-in model/server protocol, unguided-equivalence tests, noise-step critic bundle, frozen guidance configuration, matched-compute evaluation, gradient/action diagnostics, safety analysis, and videos comparing unguided, reranked, and guided executions from matched initial identities.
